% !TEX root =./ECE343_Practice_main.tex
%
% L23 practice set (Boundary Value Problems): Poisson's and Laplace's
% equations, parallel plates with and without charge, the coax and its best
% inner radius, concentric spheres, and solving by averaging.

\fullwidth{\textbf{Learning Objective 5.3: } I can calculate the electric field
from the scalar electric potential in a variety of geometries.
\\[4pt]\normalfont\small Show your work. Take $\epsilon = \epsilon_0 =
8.854\times10^{-12}$ F/m unless told otherwise. Write a \textbf{Documentation}
line at the end (\textbf{None} if you did not collaborate).}

\begin{questions}

%%%%%%%%%%%%%%%%%%%%%%%%%%%%%%%%%%%%%%%%%
\question \textbf{Poisson's and Laplace's equations.}
\begin{parts}

	\part Starting from $\nabla\cdot\mathbf{D} = \rho_v$, derive Poisson's equation
	for a homogeneous material.
	\begin{solution}[1.2in]
	\eq{
		\nabla\cdot\left(\epsilon\mathbf{E}\right) = \rho_v
		\ \Rightarrow\ \epsilon\,\nabla\cdot\left(-\nabla V\right) = \rho_v
		\ \Rightarrow\ \nabla^2 V = -\dfrac{\rho_v}{\epsilon}
	}
	$\epsilon$ leaves the divergence because it is the same everywhere.
	\end{solution}

	\part Which of these potentials can exist in a charge-free region? For any
	that cannot, find $\rho_v$ in free space.
	\begin{subparts}
		\subpart $V_1 = 3x^2 - 3y^2 + 2z$
		\subpart $V_2 = x^2 + y^2 + z^2$
		\subpart $V_3 = \ln\rho$ (cylindrical, $\rho > 0$)
		\subpart $V_4 = 1/r$ (spherical, $r > 0$)
	\end{subparts}
	\begin{solution}[1.6in]
	(1) $6 - 6 + 0 = 0$: yes.
	(2) $2 + 2 + 2 = 6$: no; $\rho_v = -6\epsilon_0 = -5.31\times10^{-11}\ \text{C/m}^3$.
	(3) $\frac{1}{\rho}\frac{d}{d\rho}\left(\rho\cdot\frac{1}{\rho}\right) = 0$: yes.
	(4) $\frac{1}{r^2}\frac{d}{dr}\left(r^2\cdot\left(-\frac{1}{r^2}\right)\right) = 0$: yes.
	\end{solution}

\end{parts}

\newpage
%%%%%%%%%%%%%%%%%%%%%%%%%%%%%%%%%%%%%%%%%
\question \textbf{Parallel plates.} Two large plates sit at $z = 0$ and $z = d$
in air.
\begin{parts}

	\begin{minipage}[t]{0.58\linewidth}
		\part The lower plate is grounded and the upper plate is at $12$ V, with
		$d = 2\ \text{mm}$. Solve Laplace's equation for $V(z)$, and find $V$ at
		$z = 0.5\ \text{mm}$ and $\mathbf{E}$.
		\begin{solution}[1.3in]
		$d^2V/dz^2 = 0$ gives $V = Az + B$; $V(0) = 0$ gives $B = 0$, and $V(d) = 12$
		gives $A = 12/d$.
		\eq{
			V(z) &= 12\,\dfrac{z}{d}, \qquad V(0.5\ \text{mm}) = 3\ \text{V} \\
			\mathbf{E} &= -\zhat\,\dfrac{12}{2\times10^{-3}} = -\zhat\,6\ \text{kV/m}
		}
		\end{solution}
	\end{minipage}\hfill%
	\begin{minipage}[t]{0.37\linewidth}
		\raggedleft \vspace{12pt}
		$V$ = \ansbox[1.3in]{$3$ V}

		\vspace{4pt}
		$\mathbf{E}$ = \ansbox[1.3in]{$-\zhat\,6$ kV/m}
	\end{minipage}

	\part Now both plates are grounded and the space between them holds a uniform
	charge density $\rho_v$. Solve Poisson's equation for $V(z)$ and find $\mathbf{E}$.
	\begin{solution}[1.6in]
	\eq{
		\dfrac{d^2V}{dz^2} &= -\dfrac{\rho_v}{\epsilon} \ \Rightarrow\ V = -\dfrac{\rho_v z^2}{2\epsilon} + Az + B \\
		V(0) = 0:\ B = 0, &\qquad V(d) = 0:\ A = \dfrac{\rho_v d}{2\epsilon} \\
		V(z) = \dfrac{\rho_v}{2\epsilon}\,z\,(d - z), &\qquad
		\mathbf{E} = -\zhat\,\dfrac{dV}{dz} = -\zhat\,\dfrac{\rho_v}{2\epsilon}\,(d - 2z)
	}
	\end{solution}

	\begin{minipage}[t]{0.58\linewidth}
		\part For $\rho_v = 1\ \mu\text{C/m}^3$ and $d = 1\ \text{cm}$, find the
		largest potential between the plates and the field at the lower plate.
		\begin{solution}[1.1in]
		\eq{
			V\!\left(\tfrac{d}{2}\right) &= \dfrac{\rho_v d^2}{8\epsilon_0} = \dfrac{10^{-6}\times10^{-4}}{8\times8.854\times10^{-12}} = 1.41\ \text{V} \\
			\mathbf{E}(0) &= -\zhat\,\dfrac{\rho_v d}{2\epsilon_0} = -\zhat\,565\ \text{V/m}
		}
		\end{solution}
	\end{minipage}\hfill%
	\begin{minipage}[t]{0.37\linewidth}
		\raggedleft \vspace{12pt}
		$V_\text{max}$ = \ansbox[1.3in]{$1.41$ V}

		\vspace{4pt}
		$\mathbf{E}(0)$ = \ansbox[1.3in]{$-\zhat\,565$ V/m}
	\end{minipage}

\end{parts}

\newpage
%%%%%%%%%%%%%%%%%%%%%%%%%%%%%%%%%%%%%%%%%
\question \textbf{The coaxial line.} An air-filled coax has inner radius $a$ and
outer radius $b$.
\begin{parts}

	\part The inner conductor is grounded and the outer conductor is at $V_0$.
	Solve Laplace's equation for $V(\rho)$ and find $\mathbf{E}$.
	\begin{solution}[1.5in]
	$\frac{1}{\rho}\frac{d}{d\rho}\left(\rho\frac{dV}{d\rho}\right) = 0$ gives
	$V = A\ln\rho + B$. $V(a) = 0$ gives $V = A\ln(\rho/a)$, and $V(b) = V_0$ gives
	$A = V_0/\ln(b/a)$:
	\eq{
		V(\rho) = V_0\,\dfrac{\ln(\rho/a)}{\ln(b/a)}, \qquad
		\mathbf{E} = -\rhohat\,\dfrac{dV}{d\rho} = -\rhohat\,\dfrac{V_0}{\rho\ln(b/a)}
	}
	The field points inward, from the outer conductor at high potential.
	\end{solution}

	\begin{minipage}[t]{0.58\linewidth}
		\part For $a = 1\ \text{mm}$, $b = 4\ \text{mm}$, and $V_0 = 100$ V, find
		$V$ at $\rho = 2\ \text{mm}$ and $\vert\mathbf{E}\vert$ at $\rho = a$.
		\begin{solution}[1.1in]
		\eq{
			V &= 100\,\dfrac{\ln 2}{\ln 4} = 50\ \text{V} \\
			\vert\mathbf{E}(a)\vert &= \dfrac{100}{10^{-3}\times1.386} = 72.1\ \text{kV/m}
		}
		\end{solution}
	\end{minipage}\hfill%
	\begin{minipage}[t]{0.37\linewidth}
		\raggedleft \vspace{12pt}
		$V$ = \ansbox[1.3in]{$50$ V}

		\vspace{4pt}
		$\vert\mathbf{E}(a)\vert$ = \ansbox[1.3in]{$72.1$ kV/m}
	\end{minipage}

	\begin{minipage}[t]{0.58\linewidth}
		\part An air coax with $b = 5\ \text{mm}$ must hold $1$ kV between its
		conductors. Find the inner radius $a$ that makes the largest field,
		$V_0/(a\ln(b/a))$, as small as possible, and that field. Air breaks down
		near $3$ MV/m; is the design safe?
		\begin{solution}[1.7in]
		Minimize $V_0/f(a)$ by maximizing $f(a) = a\ln(b/a)$:
		\eq{
			\dfrac{df}{da} &= \ln\dfrac{b}{a} - 1 = 0 \ \Rightarrow\ a = \dfrac{b}{e} = 1.84\ \text{mm} \\
			\vert\mathbf{E}\vert_\text{max} &= \dfrac{V_0}{(b/e)\times1} = \dfrac{e\,V_0}{b} = \dfrac{2.718\times1000}{0.005} = 544\ \text{kV/m}
		}
		That is well below $3$ MV/m, so the design is safe.
		\end{solution}
	\end{minipage}\hfill%
	\begin{minipage}[t]{0.37\linewidth}
		\raggedleft \vspace{12pt}
		$a$ = \ansbox[1.3in]{$1.84$ mm}

		\vspace{4pt}
		$\vert\mathbf{E}\vert_\text{max}$ = \ansbox[1.3in]{$544$ kV/m}
	\end{minipage}

\end{parts}

\newpage
%%%%%%%%%%%%%%%%%%%%%%%%%%%%%%%%%%%%%%%%%
\question \textbf{Concentric spheres.} A metal sphere of radius $a = 1\ \text{cm}$
at $V_0 = 50$ V sits at the center of a grounded spherical shell of radius
$b = 2\ \text{cm}$, with air between them.
\begin{parts}

	\part Using $\nabla^2 V = \frac{1}{r^2}\frac{d}{dr}\left(r^2\frac{dV}{dr}\right)$
	for a potential that depends on $r$ alone, solve Laplace's equation for $V(r)$
	and find $\mathbf{E}$.
	\begin{solution}[1.7in]
	$r^2\,dV/dr = A$, so $dV/dr = A/r^2$ and $V = -A/r + B$.
	$V(b) = 0$ gives $B = A/b$, so $V = -A\left(1/r - 1/b\right)$. $V(a) = V_0$ fixes $A$:
	\eq{
		V(r) = V_0\,\dfrac{1/r - 1/b}{1/a - 1/b}, \qquad
		\mathbf{E} = -\rhat\,\dfrac{dV}{dr} = \rhat\,\dfrac{V_0}{r^2\left(1/a - 1/b\right)}
	}
	\end{solution}

	\begin{minipage}[t]{0.58\linewidth}
		\part Find $V$ at $r = 1.5\ \text{cm}$ and $\vert\mathbf{E}\vert$ at the surface
		of the inner sphere.
		\begin{solution}[1.1in]
		$1/a - 1/b = 100 - 50 = 50\ \text{m}^{-1}$:
		\eq{
			V &= 50\,\dfrac{66.7 - 50}{50} = 16.7\ \text{V} \\
			\vert\mathbf{E}(a)\vert &= \dfrac{50}{10^{-4}\times50} = 10\ \text{kV/m}
		}
		\end{solution}
	\end{minipage}\hfill%
	\begin{minipage}[t]{0.37\linewidth}
		\raggedleft \vspace{12pt}
		$V$ = \ansbox[1.3in]{$16.7$ V}

		\vspace{4pt}
		$\vert\mathbf{E}(a)\vert$ = \ansbox[1.3in]{$10$ kV/m}
	\end{minipage}

\end{parts}

\vspace{1em}
%%%%%%%%%%%%%%%%%%%%%%%%%%%%%%%%%%%%%%%%%
\question \textbf{Solving by averaging.} A long square trough is divided by a grid
into four interior points, two upper ($T$) and two lower ($B$). The top wall is
at $100$ V and the other three walls are grounded. Each interior point is the
average of its four neighbors.
\begin{parts}

	\begin{minipage}[t]{0.58\linewidth}
		\part By symmetry the two upper points are equal, and so are the two lower
		points. Write the averaging rule for one upper and one lower point, and
		solve for $T$ and $B$.
		\begin{solution}[1.4in]
		An upper point has neighbors $100$, $0$ (side wall), $T$, and $B$; a lower
		point has $T$, $0$, $B$, and $0$:
		\eq{
			4T &= 100 + T + B, \qquad 4B = T + B \\
			B &= T/3, \qquad 3T = 100 + T/3 \\
			T &= 37.5\ \text{V}, \qquad B = 12.5\ \text{V}
		}
		\end{solution}
	\end{minipage}\hfill%
	\begin{minipage}[t]{0.37\linewidth}
		\raggedleft \vspace{12pt}
		$T$ = \ansbox[1.3in]{$37.5$ V}

		\vspace{4pt}
		$B$ = \ansbox[1.3in]{$12.5$ V}
	\end{minipage}

\end{parts}

\vspace{1.5em}
\noindent\textbf{Documentation:} \rule[-2pt]{0.6\linewidth}{0.4pt}

\end{questions}
